\subsection{An explicit compact-averaging bound for actual cubic support}
\label{sec:cubic-reynolds-support}

\paragraph{P-F54 (all sufficiently large actual degrees, with an explicit polynomial bound).}
For every integer $n\ge33$, there is an explicitly specified
cubic tensor $F_n\in\operatorname{Sym}^3\mathbb C^n$ such that
\[
 \begin{gathered}
 d\ge(10n)^{200},\qquad n\mid3d
 \quad\Longrightarrow\quad
 \text{some homogeneous degree-$d$ $SL_n$ invariant is nonzero at }F_n.
 \end{gathered}                                                   \tag{P-F54a}
\]
In particular, with $g=\gcd(n,3)$ and the actual cubic stretching
function $f_n$ of P-F31,
\[
 f_n(s)>0\qquad\text{for all integers }s\ge
             \left\lceil\frac{g(10n)^{200}}n\right\rceil .        \tag{P-F54b}
\]
The exponent is deliberately coarse. This is an unconditional
bound for actual support, including all allowed residue classes.
It is not a bound for the strict saturation or positivity index:
it does not evaluate a constituent polynomial away from its own
residue progression. The proof gives a compact averaging formula
for a nonzero invariant, not a claim of a polynomial-size
evaluation circuit or a positive counting formula.

\paragraph{1. A sparse exactly balanced tensor.}
Use the Hermitian tensor-product norm on
$W=\operatorname{Sym}^3\mathbb C^n\subset(\mathbb C^n)^{\otimes3}$.
The vectors $e_i^{\otimes3}$ have norm one. For a three-element
set $T$, let $v_T$ be the sum of its six ordered tensor words
divided by $\sqrt6$; the $v_T$ form an orthonormal family.
All indices in the following two lists are in $\mathbb Z/n\mathbb Z$:
\[
 E_i=\{i,i+1,i+3\},\qquad E_{n+i}=\{i,i+4,i+16\}
                      \quad(0\le i<n).                         \tag{P-F54c}
\]
These $2n$ triples are distinct, every vertex occurs in exactly six
of them, and two distinct triples intersect in at most one vertex.
Indeed their possible ordered pair differences are respectively
$\pm1,\pm2,\pm3$ and $\pm4,\pm12,\pm16$. These twelve
residues are distinct when $n\ge33$, and each difference occurs
in exactly one ordered pair within the two base triples. This
also proves that distinct translates or different lists cannot
share two vertices.

Set
\[
 \begin{gathered}
 L=100n^2,\qquad \varepsilon=n^{-20},\qquad
 \eta_j=\exp(2\pi i j/L),\quad0\le j<2n,\\
 F_0=\sum_{i=0}^{n-1}e_i^{\otimes3},\qquad
 H=\sum_{j=0}^{2n-1}\eta_jv_{E_j},\qquad
 F=F_n=F_0+\varepsilon H,\qquad
 N_F=\|F\|^2=n(1+2\varepsilon^2)\le2n.
 \end{gathered}                                                   \tag{P-F54d}
\]
Thus every coefficient is explicitly algebraic. The one-factor
partial trace of $F_0F_0^*$ is the identity matrix. That of
$v_Tv_T^*$ is $\frac13\sum_{i\in T}e_ie_i^*$.
Every cross partial trace between distinct triples is zero,
because a nonzero cross term would require two common tensor
indices. The cross partial traces between $F_0$ and $H$ are
also zero. Consequently the one-factor partial trace of $FF^*$
is exactly $(1+2\varepsilon^2)I_n$.

Let $\rho$ denote the unitary action on $W$. For a traceless
skew-Hermitian matrix $X$, write $A_X=d\rho(X)$ and
$r=\|X\|_{\mathrm F}$, where $\|\cdot\|_{\mathrm F}$ is the
matrix Frobenius norm. The partial-trace identity gives
\[
       \langle F,A_XF\rangle=0,
       \qquad \|A_X\|_{\mathrm{op}}\le3r.                    \tag{P-F54e}
\]
Here the inner product is linear in its second argument.
Directly in the orthonormal cubic monomial basis,
\[
 \|A_XF_0\|^2
   =3\sum_{i\ne j}|X_{ij}|^2+9\sum_i|X_{ii}|^2
   \ge3r^2.
\]
Since $\|H\|=\sqrt{2n}$, it follows that
\[
 \|A_XF\|\ge(\sqrt3-3\varepsilon\sqrt{2n})r\ge r
                        \qquad(n\ge33).                      \tag{P-F54f}
\]
Thus both the exact zero expectation and a quantitative
nondegeneracy bound have been proved, without a balancing
existence theorem.

\paragraph{2. Separation of the finite Fermat symmetries.}
Let $\mathcal K$ be the finite subgroup of $SU(n)$ consisting of
the monomial matrices
$h=\alpha\operatorname{diag}(\zeta_i)P$, where
$\zeta_i^3=1$, $P$ is a permutation matrix, and $\det h=1$.
Such an $h$ satisfies $\rho(h)F_0=\alpha^3F_0$.
Put $k_h=\alpha^{-3}\rho(h)$, a unitary operator fixing $F_0$.
The cube of the chosen $\alpha$ is independent of the common
cube-root ambiguity in this expression.
We claim that
\[
       h\in\mathcal K\setminus Z(SU(n))
       \quad\Longrightarrow\quad
                  \|H-k_hH\|\ge L^{-1}.                     \tag{P-F54g}
\]
If the permutation of variables takes an edge outside the
support in (P-F54c), the difference has a coefficient of
absolute value one. If it preserves the edge set but moves
one edge to another, the corresponding difference of
coefficients is $\eta_i-\zeta\eta_j$ with $i\ne j$ and
$\zeta^3=1$. For $\zeta=1$ its absolute value is at least
$4/L$, by $\sin x\ge2x/\pi$ on $[0,\pi/2]$.
For $\zeta\ne1$, it is at least one: all ratios
$\eta_i/\eta_j$ have arguments of absolute value less than
$4\pi n/L<\pi/6$. These bounds imply (P-F54g) in both cases.

If every edge is fixed individually, every vertex is fixed.
Indeed the indices of the first-list edges incident to vertex
$j$ are $\{j,j-1,j-3\}$. These sets are distinct for different
$j$. A nonzero translation preserving a three-element set
would act on it as a three-cycle, so its nonzero pair
differences would have only two values; the set
$\{0,-1,-3\}$ instead has the six distinct differences
$\pm1,\pm2,\pm3$ when $n\ge33$.
Thus the remaining case is diagonal. Write
$\zeta_i=\exp(2\pi i b_i/3)$, $b_i\in\mathbb F_3$.
If one edge product is nontrivial, the coefficient difference
has absolute value $\sqrt3$. Otherwise, for the cyclic shift
operator $S$ on the vector $(b_i)$,
\[
       (1+S+S^3)b=0,
       \qquad(1+S^4+S^{16})b=0.                              \tag{P-F54h}
\]
In $\mathbb F_3[z]$,
$1+z+z^3=(z-1)(z^2+z+2)$.
The quadratic is irreducible and its roots have fourth
power $-1$ and order eight. At either root,
$1+z^4+z^{16}=1\ne0$; at $z=1$ it vanishes.
The greatest common divisor of the two displayed polynomials
is therefore $z-1$. Bezout's polynomial identity applied to
$S$ gives $(S-1)b=0$, so all $b_i$ agree.
Then $h$ is scalar. This proves (P-F54g) in every case,
including when $3$ divides $n$.

\paragraph{3. A quantitative inverse estimate at Fermat.}
The following elementary estimate will keep the subsequent
global argument from assuming that genericity supplies a
uniform gap. Suppose $U\in SU(n)$, $|\lambda|=1$, and
\[
             \rho_0=\|F_0-\lambda\rho(U)F_0\|
                            \le\frac1{100\sqrt n}.
\]
Then there exist $h\in\mathcal K$, $X\in\mathfrak{su}(n)$,
and $t\in\mathbb R$ such that
\[
 U=h\exp X,\qquad \lambda\rho(h)=e^{it}k_h,
 \qquad \|X\|_{\mathrm F}\le24\rho_0,
 \qquad |t|\le24\rho_0/\sqrt n.                              \tag{P-F54i}
\]
To prove it, write $u_{ij}$ for the entries of $U$.
Since $\langle F_0,\rho(U)F_0\rangle=\sum_{i,j}u_{ij}^3$,
\[
 0\le n-\sum_{i,j}|u_{ij}|^3\le\rho_0^2/2.
\]
In each row let $a_i=\max_j|u_{ij}|$. The row sum of the
squares is one, so its sum of cubes is at most $a_i$.
It follows that $\sum_i(1-a_i)\le\rho_0^2/2$.
Each $a_i>1/\sqrt2$, so the maximizing positions lie in
distinct columns. They define a monomial unitary matrix
$M$, with the same phases $m_i$ at those positions, and
\[
                    \|U-M\|_{\mathrm F}^2
                       =2\sum_i(1-a_i)\le\rho_0^2.
\]
The omitted off-position cubes have total absolute value at
most $\sum_i(1-a_i^2)\le2\sum_i(1-a_i)$.
Also $\sum_i(1-a_i^3)\le3\sum_i(1-a_i)$.
Comparing the dominant-position phases with the entire cubic
sum therefore gives
\[
       \sum_i\bigl(1-\operatorname{Re}(\lambda m_i^3)\bigr)
                   \le3\rho_0^2,
       \qquad \sum_i|1-\lambda m_i^3|^2\le6\rho_0^2.
\]
Choose $\alpha^3=\lambda^{-1}$ and choose cube roots
$\zeta_i$ nearest to $m_i/\alpha$ on the unit circle.
For an angle $|\theta|\le\pi/3$,
$|e^{i\theta}-1|\le|e^{3i\theta}-1|$; hence the monomial
matrix $M'=\alpha\operatorname{diag}(\zeta_i)P$ satisfies
\[
                    \|U-M'\|_{\mathrm F}<4\rho_0.
\]
Let the principal eigenangles of $U^*M'$ be $\theta_j$.
The preceding bound and $|\theta|\le2|e^{i\theta}-1|$
on the relevant small arc imply
$|\sum_j\theta_j|\le8\sqrt n\rho_0$.
Set $\tau=n^{-1}\sum_j\theta_j$ and $h=e^{-i\tau}M'$.
Then $\det h=1$, $h\in\mathcal K$, and
$\|U-h\|_{\mathrm F}\le12\rho_0$.
The principal logarithm of $h^*U$ has Frobenius norm at
most $24\rho_0$. Its trace is a multiple of $2\pi i$,
but its absolute trace is at most
$24\sqrt n\rho_0<2\pi$, so this trace is zero.
This logarithm is the required $X$.
Finally $\lambda\alpha^3e^{-3i\tau}=e^{-3i\tau}$,
so $t=-3\tau$ gives the stated phase bound and proves
(P-F54i).

\paragraph{4. A global overlap gap.}
Put
\[
 z(U)=\frac{\langle F,\rho(U)F\rangle}{N_F},
 \qquad R=\frac{\varepsilon}{16L},
 \qquad \Delta=\frac{R^2}{16n}.
                                                               \tag{P-F54j}
\]
For every $U$ outside the union of the scalar-centre neighborhoods
\[
  \bigcup_{h\in Z(SU(n))}
       \{h\exp X:X\in\mathfrak{su}(n),\ \|X\|_{\mathrm F}<R\},
\]
we claim
\[
                            |z(U)|\le1-\Delta.               \tag{P-F54k}
\]
Choose a phase $\lambda$ minimizing
$r_F=\|F-\lambda\rho(U)F\|$.
If $r_F\ge\varepsilon/(4L)$, then $r_F\ge R/2$ already.
Otherwise
\[
 \rho_0=\|F_0-\lambda\rho(U)F_0\|
       \le r_F+2\varepsilon\sqrt{2n}
       \le5\varepsilon\sqrt n\le\frac1{100\sqrt n}.
\]
Use (P-F54i). Writing $r=\|X\|_{\mathrm F}$, we have
$r\le120\varepsilon\sqrt n$ and $|t|\le120\varepsilon$.
Project $F-\lambda\rho(U)F$ orthogonally to the squarefree
cubic subspace. The vectors $F_0$ and $A_XF_0$ have no
squarefree component, and $k_h$ preserves this subspace.
For a skew-Hermitian operator $A$,
$\|e^A-I-A\|\le\|A\|^2/2$ and
$\|e^A-I\|\le\|A\|$, by their integral remainder formulas.
Thus that projection differs from
$\varepsilon(H-k_hH)$ in norm by at most
\[
 \frac92\sqrt n\,r^2
    +\varepsilon\sqrt{2n}(|t|+3r)
 \le10^5n^{3/2}\varepsilon^2
 \le\frac{\varepsilon}{4L}.                                 \tag{P-F54l}
\]
The last inequality follows from
$4\cdot10^7 n^{7/2}\varepsilon\le1$ for $n\ge33$.
If $h$ were noncentral, (P-F54g) would give
$r_F\ge3\varepsilon/(4L)$, a contradiction.
Hence $h$ is central.

Now $k_h=I$ and the residual is
$\|F-e^{A_X+it}F\|$.
By (P-F54e)--(P-F54f),
\[
       \|(A_X+it)F\|^2
            =\|A_XF\|^2+t^2N_F\ge r^2+nt^2.
\]
For skew-Hermitian $B$, the stronger integral remainder bound
$\|(e^B-I-B)F\|\le\frac12\|B\|\,\|BF\|$ applies.
Here $\|A_X+it\|\le480\varepsilon\sqrt n<1$, so
$r_F\ge\frac12\|(A_X+it)F\|\ge r/2$.
As $U$ lies outside the displayed neighborhoods, $r\ge R$,
and therefore again $r_F\ge R/2$.
Finally
$r_F^2=2N_F(1-|z(U)|)$ and $N_F\le2n$ give
(P-F54k). All smallness inequalities used here follow directly
from $n\ge33$ and $\varepsilon=n^{-20}$; no unquantified
neighborhood or compactness gap has been used.

\paragraph{5. Local modulus and phase estimates.}
For $X\in\mathfrak{su}(n)$ with $r=\|X\|_{\mathrm F}\le R$,
Taylor's formula with the same unitary remainders gives
\[
 |z(e^X)-1|\le\tfrac92 r^2,
 \qquad |\operatorname{Im}z(e^X)|\le\tfrac92 r^3,
 \qquad |\arg z(e^X)|\le10r^3.                               \tag{P-F54m}
\]
The linear term is zero by (P-F54e); the quadratic term
$\langle F,A_X^2F\rangle=-\|A_XF\|^2$ is real.
The second and third Taylor remainder bounds are respectively
$\frac12\|A_X\|^2$ and $\frac16\|A_X\|^3$ after dividing
by $N_F$. Since the real part is at least $1/2$, these also
give the argument bound.

Taking $t=-\arg z(e^X)$ and using the local residual argument
from step~4 gives
\[
                         |z(e^X)|\le1-\frac{r^2}{16n}.
                                                               \tag{P-F54n}
\]
Indeed $\|A_X+it\|\le3r+10r^3<1$; the exact zero
expectation removes the mixed term in its norm, and
(P-F54f) gives residual at least $r/2$.

\paragraph{6. A Haar-volume lower bound without a dimension-free shortcut.}
Normalize Haar measure $\mu$ on $SU(n)$ to total mass one.
For $0<a\le1/(10\sqrt n)$,
\[
 \mu\{e^X:X\in\mathfrak{su}(n),\ \|X\|_{\mathrm F}<a\}
                   \ge(1+8\sqrt n/a)^{-2n^2}.                \tag{P-F54o}
\]
Here is a direct covering proof. A maximal collection of
points of $SU(n)$ separated by Frobenius distance $b$ has
at most $(1+2\sqrt n/b)^{2n^2}$ members: the disjoint
ambient Euclidean balls of radius $b/2$ around them lie
in a ball of radius $\sqrt n+b/2$ in
$\mathbb R^{2n^2}$. The group balls of radius $b$ cover,
and their Haar measures agree by unitary invariance of
Frobenius distance. The identity ball therefore has
measure at least the reciprocal of this number.
Take $b=a/4$, giving the bound
$(1+8\sqrt n/a)^{2n^2}$ exactly.
Every unitary in this identity ball has a principal logarithm
of norm at most $2b=a/2<a$. Its trace is a multiple of $2\pi i$
and has absolute value at most $\sqrt n a<2\pi$, so the
logarithm belongs to $\mathfrak{su}(n)$.
The closed identity ball of radius $b$ is thus contained
in the open logarithm ball of radius $a$, proving (P-F54o)
without an assertion about boundary measure.

\paragraph{7. The complete compact-average comparison.}
Let $d$ be an integer with $d\ge(10n)^{200}$ and $n\mid3d$.
Set
\[
             b_d=d^{-2/5},\qquad a_d=\frac1{10\sqrt d}.
\]
We have $a_d\le b_d\le R$.
For $r\le b_d$, (P-F54m) gives
$d|\arg z(e^X)|\le10d^{-1/5}\le1$.
Thus the real part of $z(e^X)^d$ is nonnegative on this
whole neighborhood. On the smaller neighborhood $r\le a_d$,
\[
 |z(e^X)|\ge1-\frac1{20d},
 \qquad \operatorname{Re}z(e^X)^d\ge\frac14.                \tag{P-F54p}
\]
For the latter use $(1-1/(20d))^d\ge19/20$ and
$\cos(d\arg z)\ge1/2$.

A scalar center element $h=\omega I_n$, $\omega^n=1$,
acts on $W$ by $\omega^3$. Hence
$z(he^X)^d=\omega^{3d}z(e^X)^d=z(e^X)^d$.
All central neighborhoods therefore have the same positive
phase; none are discarded by an absolute-value argument.
Outside their union of radius $b_d$, either (P-F54k) or
(P-F54n) applies and gives
\[
       |z(U)|^d\le
           \exp\left(-\frac{d^{1/5}}{16n}\right).             \tag{P-F54q}
\]
The neighborhoods can be taken disjoint: distinct scalar
centers have Frobenius distance at least $4/\sqrt n$,
whereas $2R<4/\sqrt n$. Combining (P-F54o)--(P-F54q)
and retaining just the smaller identity neighborhood yields
\[
 \operatorname{Re}\int_{SU(n)}z(U)^d\,d\mu(U)
 \ge\frac14(1+80\sqrt{nd})^{-2n^2}
        -\exp\left(-\frac{d^{1/5}}{16n}\right)>0.             \tag{P-F54r}
\]
For clarity, the final strict inequality follows from explicit
elementary bounds. Since $nd\ge1$,
$1+80\sqrt{nd}\le100nd$.
When $d\ge(10n)^{200}$,
$\log(100n)\le\frac1{100}\log d$, so
\[
 \log4+2n^2\log(1+80\sqrt{nd})<5n^2\log d.
\]
Also $d^{1/5}>80n^3\log d$ on this range.
To check the latter at the left endpoint, substitute
$d=(10n)^{200}$ and use $\log(10n)\le10n$;
$10^{40}n^{40}>160000n^4$ suffices.
The function $d^{1/5}/\log d$ is increasing for $d>e^5$,
so the inequality persists for every larger $d$.
This proves the strict positivity in (P-F54r), with the
full dependence on $n$ in the neighborhood volumes included.

\paragraph{8. From the average to an actual polynomial invariant.}
The compact averaging operation is the usual Reynolds
construction; an exact modern primary locator is Aryaman
Maithani, \emph{Splitting the difference: Computations of the
Reynolds operator in classical invariant theory},
\emph{Bulletin of the London Mathematical Society} 58 (2026),
e70175, Theorem~5.6. Theorem~6.1(a) gives the density of
$SU(n)$ in $SL_n(\mathbb C)$:
\url{https://doi.org/10.1112/blms.70175}.
The following short argument records the precise tensor-norm
identity needed for the quantitative estimate above; no degree
bound is being attributed to those general statements.

The Haar average $P_d=\int_{SU(n)}\rho(U)^{\otimes d}d\mu(U)$
is the orthogonal projection onto the invariant subspace:
left invariance makes its image invariant, it fixes every
invariant vector, and inversion invariance makes it
self-adjoint. Therefore
\[
 \int_{SU(n)}z(U)^d\,d\mu(U)
             =\frac{\|P_dF^{\otimes d}\|^2}{N_F^d}.           \tag{P-F54s}
\]
The tensor $F^{\otimes d}$ and its projection lie in
$\operatorname{Sym}^dW$.
By (P-F54r) the projection is nonzero.
The polynomial
$I(v)=\langle P_dF^{\otimes d},v^{\otimes d}\rangle$
is homogeneous of degree $d$, invariant under $SU(n)$,
and satisfies $I(F)=\|P_dF^{\otimes d}\|^2>0$.
It is invariant under $SL_n(\mathbb C)$ as well:
the Lie algebra $\mathfrak{su}(n)$ annihilates it, its
complex span is $\mathfrak{sl}_n(\mathbb C)$, and the
elementary unipotent one-parameter subgroups generate
$SL_n(\mathbb C)$. This proves (P-F54a), and substituting
$d=ns/g$ proves (P-F54b).

The center condition is also necessary: the scalar center
acts on degree-$d$ cubic tensors by $\omega^{3d}$.
Thus this one explicitly specified point has every
sufficiently large degree permitted by the minimal center
condition. No assertion about $n<33$ is needed for the
uniform asymptotic statement, and no generic stabilizer
theorem, GRH, numerical semigroup conductor theorem, or
unquantified stationary-phase limit enters the proof.
The distinction from strict constituent positivity remains
exactly the one stated at the beginning.\hfill$\square$

\paragraph{An exact illustration of the remaining index distinction.}
For any integer $M\ge2$, consider the quasipolynomial
\[
 q_M(s)=M\lfloor s/M\rfloor+1,
 \qquad Q_i(X)=X-i+1\quad(0\le i<M).                       \tag{P-F54t}
\]
It is positive at every actual integer $s\ge0$, and
$Q_i(s)=q_M(s)$ when $s\equiv i\pmod M$.
The bounded periodic function $q_M(s)-s=1-(s\bmod M)$
has minimal period $M$, so these are its constituents.
A shift by a nonnegative integer $c$ permutes the
constituents $Q_i(X+c)$; their smallest constant
coefficient is $c-M+2$. Hence all coefficients are
nonnegative exactly when $c\ge M-2$.
For strict saturation, each constituent is increasing
in $X$, and its smallest value at a positive integer
is $Q_{M-1}(1+c)=c-M+3$. This is strictly positive
exactly when $c\ge M-2$. Thus the literal conventions
of GCT VI, Definitions~1.2.2 and 1.2.4, printed
pp.~13--14, give
\[
                   p(q_M)=s(q_M)=M-2.                    \tag{P-F54u}
\]
See \url{https://arxiv.org/pdf/0704.0229}.
This sequence is not asserted to be a plethysm sequence.
It only demonstrates why even positivity at every actual
nonnegative integer does not by itself bound the two
strict constituent indices.

\paragraph{The same distinction for a normal Gorenstein quotient surface.}
The distinction persists for actual Hilbert functions with much of
the geometry used in the proposed purity bridge. Let $M\ge2$ be even
and let $\mu_{2M}$ act on $\mathbb C[u,v]$ by
$(u,v)\mapsto(\zeta u,\zeta^{-1}v)$. Regrade its invariant ring
by half the ordinary polynomial degree. There is an explicit graded
isomorphism
\[
 \begin{gathered}
 A_M=\mathbb C[u,v]^{\mu_{2M}}
   \cong\mathbb C[x,y,z]/(yz-x^{2M}),\\
 (x,y,z)=(uv,u^{2M},v^{2M}),\qquad
 (\deg x,\deg y,\deg z)=(1,M,M).
 \end{gathered}                                         \tag{P-F54v}
\]
Indeed an invariant monomial $u^av^b$ has $a-b$ divisible
by $2M$. If $a\ge b$ it is $x^by^{(a-b)/(2M)}$, and the
opposite inequality gives a monomial in $x,z$.
The relation displayed in (P-F54v) reduces every monomial
to one of
$x^i$, $x^iy^j$, or $x^iz^j$ with $i\ge0,j\ge1$.
Their images are distinct monomials in $u,v$, proving that
there are no further relations.

The ring is the toric semigroup ring of
\[
 \mathbb R_{\ge0}^2\cap
   \{(a,b)\in\mathbb Z^2:a-b\equiv0\pmod {2M}\}.
\]
Its normality and rational singularities are standard toric
facts, with exact primary locators Danilov,
\emph{The geometry of toric varieties}, Proposition~3.2,
printed p.~106, and Proposition~8.5.1, pp.~124--125:
\url{https://www.math.purdue.edu/~arapura/danilov.pdf}.
It is a two-dimensional hypersurface, hence Gorenstein
and Cohen--Macaulay; see Stacks Project,
Lemma~47.21.8, Tag~\texttt{0DWA},
\url{https://stacks.math.columbia.edu/tag/0DWA}.
Its $a$-invariant is $2M-(1+M+M)=-1$.
Equivalently, applying graded $\operatorname{Hom}$ to the
one-relation free resolution gives its canonical module
$A_M(2M-(1+2M))=A_M(-1)$.

The monomial basis just proved gives, for every $s\ge0$,
\[
 \operatorname{Hilb}_{A_M}(t)
    =\frac{1+t^M}{(1-t)(1-t^M)},\qquad
 h_M(s)=\dim(A_M)_s=2\lfloor s/M\rfloor+1>0.              \tag{P-F54w}
\]
Thus its constituents are
$H_i(X)=2X/M+1-2i/M$, $0\le i<M$.
After a shift $c\ge0$, the smallest constant coefficient
is $(2c-M+2)/M$, and the smallest value at a positive
integer argument is $(2c-M+4)/M$.
The first is nonnegative exactly when $c\ge M/2-1$;
the second is strictly positive exactly when $c\ge M/2-1$,
because $c$ and $M/2$ are integers. Consequently
\[
             p(h_M)=s(h_M)=M/2-1 \qquad(M\ge2\text{ even}).
                                                               \tag{P-F54x}
\]
This is a normal Gorenstein rational surface of fixed
dimension and fixed $a$-invariant, with every actual degree
present and arbitrarily large strict indices. It is not
asserted to be an ordinary plethysm algebra.
Its positive rational form also contains a denominator
factor $1-t$. Thus that factor alone does not imply the
literal strict saturation of Definition~1.2.4 when the
nonnegative Hilbert numerator is unrestricted. The sentence
after GCT VI, Definition~1.2.7, requires this qualification;
the conflicting printed numerator-degree index was already
audited in P-F12. None of these scope observations supplies
an upper or lower bound for the actual cubic strict indices.

\paragraph{A nonexplicit extension to every $n\ge4$.}
There is an absolute constant $C$ such that
\[
                f_n(s)>0\quad
       \text{for all }n\ge4\text{ and all integers }s\ge Cn^{199}.
                                                               \tag{P-F54y}
\]
This adds no explicit value of $C$ for the finitely many
dimensions $4\le n\le32$.
For those dimensions, B\"urgisser--Ikenmeyer,
\emph{Fundamental invariants of orbit closures},
Theorem~2.3 p.~5, Proposition~2.10 p.~7, Theorem~3.4 p.~8,
and Definition~3.6 with the following paragraph p.~9,
show that the group generated by ordinary invariant
degrees of cubic forms is $(n/\gcd(n,3))\mathbb Z$:
\url{https://arxiv.org/pdf/1511.02927}.
Hence the canonical stretch support has greatest common
divisor one, exactly as used in P-F32.
Choose finitely many nonzero homogeneous algebra generators.
Their degree semigroup is the full support: their products
are nonzero because the invariant ring is a domain, and
every homogeneous element is a polynomial in them.
For any positive semigroup member $a$, its residues form
a subgroup of $\mathbb Z/a\mathbb Z$, which is the whole
group by the greatest-common-divisor condition.
Choose one semigroup representative of each residue;
adding multiples of $a$ proves that all sufficiently
large stretch degrees occur. Let $c_n$ be one such
conductor for each $4\le n\le32$.
Taking
\[
 C\ge\max\left\{3\cdot10^{200}+1,
             \max_{4\le n\le32}c_n/n^{199}\right\}
\]
proves (P-F54y), since (P-F54b) already handles every
$n\ge33$. This finite absorption is qualitative; the
explicit balanced point and numerical threshold remain
those of (P-F54a)--(P-F54b) for $n\ge33$.
